\begin{frame}{자주 묻는 질문 (1/2)}
  \begin{itemize}
    \item \kq{Q. ``TD(0)의 (0)''은?}
      한 스텝만 앞을 내다보고 바로 부트스트래핑(0-step lookahead 이후
      부트스트랩). Week 7의 n-step TD는 $n$ 스텝 실제 관찰 후
      부트스트랩: $n = 1$이면 TD(0), $n = \infty$(끝까지)면 사실상 MC.
    \item \kq{Q. 부정확한 $V(s')$를 쓰는데도 왜 정확히 수렴하나?}
      처음엔 $V(s')$이 부정확하나, 매 갱신마다 그 오차가 TD 오차만큼
      줄고, 줄어든 값이 다시 다음 목표값에 쓰이면서 오차가 \textbf{전체적으로}
      계속 줄어든다 --- Week 4.3 바나흐 고정점과 같은 원리(차이: 샘플 기반).
    \item \kq{Q. 에피소드 내 갱신 순서(in-place)가 결과에 영향?}
      단기적으로는 조금, \textbf{수렴 결과(고정점)에는 영향 없음}.
      synchronous와 in-place는 둘 다 같은 고정점 $V^\pi$로 수렴 ---
      \textbf{일관되게} 한 방식을 쓰면 된다.
  \end{itemize}
\end{frame}
